% Cover sheet for the combined specimen PDF.
\documentclass[11pt]{article}
\usepackage[paperwidth=5.5in,paperheight=8.5in,textwidth=4.1in,top=0.8in,bottom=0.8in]{geometry}
\usepackage{fontspec}
\setmainfont{EB Garamond}[Numbers={OldStyle,Proportional}]
\usepackage{microtype}
\usepackage{xcolor}
\definecolor{grapeink}{HTML}{5B2A4E}
\pagestyle{empty}
\setlength\parindent{0pt}\setlength\parskip{0.5\baselineskip}
\begin{document}
\begin{center}
{\Huge Grapes}\par\medskip
{\scshape\color{grapeink} layout specimens}\par
\end{center}
\bigskip
The same passage, or part of it, is set five ways. Each specimen tests a way
to carry notes on notes.

\textbf{1. Stacked footnotes.} One column. Numbered notes at the foot; notes
on those notes, lettered, below them; notes on those, marked with symbols and
run together, at the very bottom.

\textbf{2. Sidenotes, their notes at the foot.} Numbered notes in a wide
outer margin. Notes on a sidenote drop to the foot of the page.

\textbf{3. Sidenotes, their notes in the margin.} As 2, but the notes on a
sidenote stay under it in the margin, in smaller type.

\textbf{4. A Talmud-style page.} A short main text at the top centre, two
commentaries wrapped around it in two typefaces, and a band of notes on the
commentaries at the foot. Composed by hand, one page at a time.

\textbf{5. Notes keyed to line numbers.} The layout of a scholarly edition.
The text carries no marks; three series of notes at the foot cite the line
number and repeat the words they comment on.

\vfill
{\small\color{grapeink} The prose is placeholder. Its facts come from memory
and have not been checked against sources.}
\end{document}
